% from ega4-18-fr.tex
\oldpage[IV-4]{195}

conclut que $x$ et $y$ sont inversibles, donc il en est de même de $(x,y)$, ce qui achève de prouver (ii).
  \item[(iii)] L'image de $\mathfrak m$ par l'homomorphisme surjectif $G\to E$ est l'idéal maximal $\mathfrak r$ de $E$, donc l'image de $\mathfrak m^n\cap\mathfrak J'$ est $\mathfrak r^n\cap\mathfrak J$, et comme $\mathfrak J$ est fermé (donc complet) pour la topologie induite par la topologie $\mathfrak r$-adique de $E$ ($0_{\mathrm I}$, 7.3.5), $\mathfrak J'$ est complet pour la topologie induite par la topologie $\mathfrak m$-préadique de $G$. D'autre part, $G/\mathfrak J'$, muni de la topologie $\mathfrak m$-préadique, est isomorphe à $F$ muni de la topologie $\mathfrak s$-adique, donc est complet pour la topologie quotient de la topologie $\mathfrak m$-préadique de $G$. On en déduit aussitôt que $G$ est complet pour la topologie $\mathfrak m$-préadique.
\end{enumerate}

Ce lemme étant établi, il résulte du théorème de Cohen ($0$, 19.8.8) que $B''$ est isomorphe à un quotient d'un anneau local noethérien complet \emph{régulier} $C$. Il résulte alors de (19.1.2) que l'immersion $\Spec(B')\to\Spec(C)$ est régulière; comme par hypothèse il en est de même de l'immersion $\Spec(A)\to\Spec(B')$, on conclut de (19.1.5, (iii)) que l'immersion $\Spec(A)\to\Spec(C)$ est régulière; mais cette immersion s'écrit aussi comme composée $\Spec(A)\to\Spec(B)\to\Spec(C)$. Or comme $B$ est régulier par hypothèse, l'immersion $\Spec(B)\to\Spec(C)$ est régulière (19.1.2); donc il en est de même de $\Spec(A)\to\Spec(B)$ par (19.1.5, (iv)).

\begin{corollary}[19.3.3]
\label{IV.19.3.3}
Soit $X$ un préschéma localement noethérien localement immersible dans un préschéma régulier. Alors l'ensemble des $x\in X$ où $X$ est une intersection complète est ouvert dans $X$ et contient l'ensemble des points réguliers de $X$.
\end{corollary}

On peut en effet se borner au cas où $X=\Spec(B/\mathfrak J)$, où $B$ est un anneau régulier et $\mathfrak J$ un idéal de $B$. L'ensemble des $\mathfrak p\supset\mathfrak J$ dans $\Spec(B)$ tels que $\mathfrak J_{\mathfrak p}$ soit régulier dans $B_{\mathfrak p}$ est alors ouvert (16.9.5), et c'est l'ensemble des points de $\Spec(B)$ où $X$ est une intersection complète en vertu de (19.3.2).

\begin{corollary}[19.3.4]
\label{IV.19.3.4}
Soient $k$ un corps, $X$ un préschéma localement de type fini sur $k$, $k'$ une extension de $k$, $X'=X\otimes_k k'$. Soient $x'$ un point de $X'$, $x$ sa projection dans $X$. Pour que $X$ soit intersection complète en $x$, il faut et il suffit que $X'$ soit intersection complète en $x'$.
\end{corollary}

La question étant locale sur $X$, on peut supposer que $X=\Spec(A)$, où $A$ est une $k$-algèbre de type fini, donc quotient d'une algèbre de polynômes $k[T_1,\ldots,T_n]$, si bien que $X$ est sous-préschéma fermé du schéma régulier $Y=\Spec(k[T_1,\ldots,T_n])$ ($0$, 17.3.7). Si l'on pose $Y'=Y\otimes_k k'=\Spec(k'[T_1,\ldots,T_n])$, $Y'$ est aussi un schéma régulier. Or, puisque $\Spec(k')$ est fidèlement plat sur $\Spec(k)$, il résulte de (19.1.5, (ii)) et du fait qu'il s'agit de préschémas noethériens que l'immersion $Y\to X$ est régulière au point $x$ si et seulement si l'immersion $Y'\to X'$ l'est au point $x'$. La conclusion résulte donc du critère (19.3.2).

\begin{corollary}[19.3.5]
\label{IV.19.3.5}
Soient $A$, $C$ deux anneaux locaux noethériens, $\rho:A\to C$ un homomorphisme local, $B=C/\mathfrak J$ un anneau quotient de $C$, $k$ le corps résiduel de $A$. On suppose que $\rho$ fait de $C$ un $A$-module plat, et que l'anneau $C\otimes_A k$ est régulier. Alors les conditions suivantes sont équivalentes:
\begin{enumerate}
  \item[a)] L'idéal $\mathfrak J$ est transversalement régulier relativement à $A$.
  \item[b)] $B$ est un $A$-module plat, et $B\otimes_A k$ est un anneau d'intersection complète.
\end{enumerate}
\end{corollary}

